\documentclass[../../../include/open-logic-section]{subfiles}

\begin{document}

\olfileid[hi]{his}{set}{mythology}

\olsection{इतिहास से अधिक मिथक?}

एक शताब्दी से अधिक समय बीतने के बाद इन घटनाओं को देखते हुए हमें सावधान रहना
चाहिए कि इन निर्णयों को केवल विश्वास के आधार पर न मान लें। परिणाम निश्चय ही
नए, रोमांचक और आश्चर्यजनक थे। किंतु वे सच में कितने चौंकाने वाले थे? और क्या
उन्होंने वास्तव में दिखाया कि हमें ज्यामितीय अंतर्ज्ञान पर भरोसा नहीं करना
चाहिए?

चकित होने के प्रश्न पर \citet{Gouvea2011} बताते हैं कि डेडेकिंड को कांटोर की
प्रसिद्ध टिप्पणी, ``\emph{je le vois, mais je ne le crois pas}'', को संदर्भ
से काफ़ी अलग करके उद्धृत किया जाता है। उसका कुछ और संदर्भ यहाँ है
(\citeauthor{Gouvea2011} से उद्धृत):
\begin{quote}
यदि मैं आपकी उदारता और धैर्य से इतनी अधिक अपेक्षाएँ कर रहा हूँ, तो विषय के
प्रति मेरे उत्साह को क्षमा करें; हाल में जो बातें मैंने आपको भेजीं, वे मेरे
लिए भी इतनी अप्रत्याशित, इतनी नई हैं कि जब तक आप, आदरणीय मित्र, उनकी शुद्धता
पर निर्णय नहीं देते, मुझे मन की शांति नहीं मिल सकती। जब तक आप मुझसे सहमत
नहीं हुए हैं, मैं केवल यही कह सकता हूँ: \emph{je le vois, mais je ne le
crois pas.}
\end{quote}
कांटोर जानते थे कि उनका परिणाम ``इतना अप्रत्याशित, इतना नया'' था। किंतु यह
संदिग्ध है कि उन्होंने कभी अपने परिणाम को \emph{अविश्वसनीय} माना हो। जैसा
\citeauthor{Gouvea2011} बताते हैं, वे केवल डेडेकिंड से अपना दिया हुआ प्रमाण
जाँचने को कह रहे थे।

ज्यामितीय अंतर्ज्ञान के प्रश्न पर: पियानो ने अपना अंतरिक्ष-पूर्ति वक्र बिना
किसी आरेख के प्रकाशित किया। किंतु जब हिल्बर्ट ने अपना वक्र प्रकाशित किया, तो
उन्होंने उद्देश्य समझाया: यदि पाठक ``निम्नलिखित ज्यामितीय अंतर्ज्ञान की सहायता
लें'', तो वे उन्हें पियानो के परिणाम को समझने का स्पष्ट ढंग देंगे; इसके बाद
उन्होंने ठीक \olref[his][set][pathology]{sec} में दिए गए आरेखों जैसे
\emph{आरेखों} की एक शृंखला सम्मिलित की।

अधिक सामान्यतः: यद्यपि प्रकाशित प्रमाणों में आरेख कुछ अप्रचलित हो गए हैं, इस
तथ्य से बचा नहीं जा सकता कि गणितज्ञ बातें सिद्ध करते समय \emph{बहुधा} आरेखों
का उपयोग करते हैं। (मोटे तौर पर: अच्छे गणितज्ञ जानते हैं कि वे ज्यामितीय
अंतर्ज्ञान पर कब भरोसा कर सकते हैं।)

संक्षेप में: प्रचार पर विश्वास न करें; या कम-से-कम उसे केवल विश्वास के आधार
पर न मानें। इस विषय में अधिक जानने के लिए \citet{Giaquinto2007} पढ़ सकते
हैं।

\end{document}
